% declaracion-ia.tex — Declaración de Uso de IA.
% main.tex la incluye DESPUÉS de Referencias, como exige la orientación
% del seminario (secc. 3): epígrafe independiente, al final del
% documento, que indique (1) herramienta y versión, (2) propósito y
% (3) forma de integración del resultado.

\section*{Declaración de Uso de la Inteligencia Artificial}
\addcontentsline{toc}{section}{Declaración de Uso de la Inteligencia Artificial}

Para este seminario, el equipo utilizó varias herramientas de inteligencia
artificial. Heily Rodríguez Rodríguez (P1 y P4) y Ernesto Alejandro Soler
Choong (P2) usaron Claude, de Anthropic, a través de la aplicación Claude Code, con
los modelos Claude Sonnet~5 y Claude Opus~5.5 respectivamente. Enrique
González Moreira (P3 y P6) usó Gemini, de Google, y GLM~5.3, de Z.ai.
Mauricio Bridon (P5) usó Codex, de OpenAI.

Las herramientas se emplearon para localizar en la bibliografía los
capítulos y pasajes pertinentes para cada inciso, redactar los borradores de
las respuestas en Markdown y en LaTeX, revisar que el estilo de citas fuera
consistente en formato APA, y compilar el documento para comprobar el
formato. En el caso de P2, también se usaron para recalcular las cifras de la
fórmula de colisiones de cachés replicadas.

Los borradores generados se utilizaron como base de las respuestas, que cada
integrante revisó, corrigió y aprobó hasta su versión final. Las citas, los
números de capítulo y los cálculos se cotejaron con los textos originales de
la bibliografía.

Todo el contenido generado fue revisado, verificado y adaptado por el
equipo, quien asume la responsabilidad final del documento entregado.
